% =============================================================================
% European Banks — Sector Intel Pack
% Demonstration artifact produced via Claude Cowork (personal instance)
% using the Market Researcher agent / sector-overview skill.
% Compile: latexmk -pdf -xelatex european-banks-sector-pack.tex
% =============================================================================
\documentclass[10pt,a4paper]{article}

\usepackage[a4paper,margin=18mm,top=15mm,bottom=18mm]{geometry}
\usepackage{parskip}
\setlength{\parskip}{4pt plus 1pt}

\usepackage[english]{babel}
\usepackage{microtype}
\usepackage{fontspec}
\usepackage{amssymb}

\usepackage{xcolor}
\definecolor{accentpurple}{HTML}{7C3AED}
\definecolor{darkgray}{HTML}{1F2937}
\definecolor{lightgray}{HTML}{F3F4F6}

\usepackage[most]{tcolorbox}

\usepackage{booktabs}
\usepackage{array}
\usepackage{tabularx}

\usepackage{titlesec}
\titleformat{\section}
  {\color{accentpurple}\bfseries\large}
  {\thesection.}{6pt}{}
  [\vspace{-4pt}\noindent\color{accentpurple}\rule{\textwidth}{0.6pt}]
\titleformat{\subsection}
  {\color{darkgray}\bfseries\normalsize}{}{0pt}{}
\titlespacing*{\section}{0pt}{10pt}{4pt}
\titlespacing*{\subsection}{0pt}{6pt}{2pt}

\usepackage{enumitem}
\setlist{nosep, leftmargin=*, topsep=2pt, itemsep=2pt}

\usepackage{hyperref}
\hypersetup{
  colorlinks=true,
  linkcolor=accentpurple,
  urlcolor=accentpurple,
  pdftitle={European Banks — Sector Intel Pack (demo artifact)}
}

\pagestyle{plain}

% =============================================================================
\begin{document}

\noindent
\begin{tcolorbox}[colback=darkgray, colframe=darkgray, arc=2pt, boxrule=0pt]
  \color{white}\textbf{\Large European Banks --- Sector Intel Pack}\\[2pt]
  \small \textbf{Prepared:} 8 June 2026 \quad \textbar\quad \textbf{Skill:} \texttt{sector-overview} (market-researcher) \quad \textbar\quad \textbf{Surface:} Claude Cowork (personal instance)
\end{tcolorbox}

\textbf{Scope:} Eurozone-listed G-SIBs and large national champions --- BNP Paribas, Banco Santander, ING Group, UniCredit, Intesa Sanpaolo, Deutsche Bank, BBVA, Soci\'et\'e G\'en\'erale.\\
\textbf{Excluded:} pure-play asset managers, pure-play investment banks, UK-domestic banks, fintech challengers.\\
\textbf{``Recent'' window:} Q1 2026 earnings season (results released late April--May 2026).

\begin{tcolorbox}[colback=lightgray, colframe=accentpurple, boxrule=1pt, leftrule=3pt, arc=2pt]
\textbf{Caveat --- read first.} This pack was assembled from public web sources (company press releases, regulator publications, market aggregators) in a personal Cowork instance with no S\&P / FactSet / LSEG connector. Market-data points (price, market cap, P/TBV) come from third-party aggregators dated late-May to 7-June 2026 and are approximate. Every numeric claim carries a source; items that could not be verified are listed explicitly in the final section. Sector packs age fast --- figures are as-of early June 2026.
\end{tcolorbox}

% =============================================================================
\section{Sector map}

The eight names span the full spread of the European banking business model, from emerging-market retail compounders to deep-discount wholesale restructurings. Three axes organise the group.

\textbf{By business mix.} At one end sit the diversified retail-and-commercial franchises whose earnings are geared to net interest income and loan growth: Santander, BBVA and ING. At the other end sit the fee-and-capital-light models that lean on wealth management, insurance and trading to offset rate pressure: Intesa Sanpaolo (its Wealth Management, Protection \& Advisory engine), UniCredit (fees, wealth and trading), and Deutsche Bank (CIB plus DWS asset management). BNP Paribas and Soci\'et\'e G\'en\'erale are universal banks straddling both --- domestic retail plus a large corporate-and-investment bank, with BNP's scaled asset-management arm (post-AXA IM) a growing third leg.

\textbf{By geography.} Santander and BBVA are the outliers, with the majority of profit sourced outside the Eurozone --- Brazil, Mexico, the US and the UK for Santander; Mexico, Turkey and South America for BBVA. That makes both structurally less exposed to ECB rate cuts but more exposed to EM FX and cost-of-risk. The rest are Eurozone-centred: BNP (France/Belgium/Italy/Luxembourg), SocGen (France plus CEE/Africa and the Ayvens mobility arm), Deutsche (Germany-anchored global CIB), ING (Benelux core plus ING-DiBa in Germany), and the two Italians (Intesa domestic; UniCredit pan-European across Italy, Germany's HVB, Austria and CEE).

\textbf{By balance-sheet shape and valuation.} The sector re-rated violently through 2025 --- the STOXX Europe 600 Banks index rose \textbf{+67\%, its best year since 1987} (\href{https://stoxx.com/european-bank-stocks-have-record-year-lifting-sector-stoxx-etf-assets-above-eur-13bn/}{STOXX}) --- but the names did not move together. ING ($\sim$1.48x P/TBV) and UniCredit ($\sim$1.88x) trade at clear premiums to tangible book on the back of sector-leading returns; BNP ($\sim$0.9--1.0x), Deutsche ($\sim$0.77x) and Soci\'et\'e G\'en\'erale ($\sim$0.6--0.8x) still trade at or below book, the latter the deepest-discount restructuring story in the group (\href{https://www.gurufocus.com/term/pb/XAMS:INGA}{gurufocus}, \href{https://companiesmarketcap.com/deutsche-bank/}{companiesmarketcap}). Against history the sector as a whole reached a euro-area aggregate price-to-book not seen since before the GFC by February 2026 before pulling back (\href{https://www.ecb.europa.eu/press/financial-stability-publications/fsr/focus/2026/html/ecb.fsrbox202605_03~ae5ac7a1de.en.html}{ECB FSR, May 2026}). European banks still trade at roughly \textbf{9.2x forward earnings versus $\sim$12.2x for US banks}, a gap that has narrowed as EU profitability converged toward US levels (EU bank ROE averaged $\sim$10\% in 2025).

The macro backdrop sets up the central debate for the year. The ECB cut its deposit rate from a 4.00\% peak to \textbf{2.00\% by June 2025 ($-$200bps cumulative)} and has held there through its March 2026 meeting (\href{https://www.ecb.europa.eu/stats/policy_and_exchange_rates/key_ecb_interest_rates/html/index.en.html}{ECB key rates}, \href{https://www.ecb.europa.eu/press/pr/date/2026/html/ecb.mp260319~3057739775.en.html}{ECB Mar 2026 decision}). The easing cycle that pressured NII through 2025 now appears finished; into the 11 June 2026 meeting, market commentary has flipped toward a possible first \textbf{hike} on inflation and geopolitical (Middle East war) risk (\href{https://www.cnbc.com/2026/04/30/european-central-bank-april-2026-rate-decision-inflation-stagflation-risk-iran-war.html}{CNBC, Apr 2026}). A steepening curve is a tailwind for the NII-geared names heading into H2.

% =============================================================================
\section{Issuer cards}

Figures are Q1 2026 unless noted. Capital return references include the most recent declared dividend and any 2025--26 buyback. Valuation is approximate early-June 2026 market data.

\subsection*{BNP Paribas (BNP.PA)}
The EU's largest bank by assets; universal model across commercial banking, CIB, and Investment \& Protection Services (now housing the integrated AXA IM). Q1 net income \textbf{\texteuro3,217m (+9\% YoY)} on revenue \textbf{\texteuro14,056m (+8.5\%)}, with IPS up $\sim$33\% on the AXA IM consolidation (\href{https://invest.bnpparibas/en/document/1q26-pr}{BNP 1Q26}). CET1 \textbf{12.8\%}, target raised to 13\% by 2027; cost of risk \textbf{39bps}; FY26 RoTE target \textbf{12\%}, rising to 13\% by 2028. Capital return: 2025 dividend \texteuro5.16/share plus a \texteuro1.15bn buyback completed Dec 2025. Trades $\sim$0.9--1.0x P/TBV, $\sim$8.8x P/E, market cap $\sim$\texteuro102bn. \textbf{Thesis driver:} commercial-bank NII tailwind from a steepening curve plus the AXA IM build-out, with a credible path to a 13\% RoTE.

\subsection*{Banco Santander (SAN)}
Globally diversified retail/commercial bank --- Spain, Brazil, Mexico, UK and a scaling US franchise. Q1 underlying net profit a record \textbf{\texteuro3.56bn (+12\%)}; NII \textbf{\texteuro11.02bn (+4\%)}; underlying RoTE \textbf{15.2\%}, targeting \textbf{>20\% by 2028} (\href{https://www.santander.com/en/press-room/press-releases/2026/04/q1-2026-santander-bank-results}{Santander Q1 2026}). CET1 \textbf{14.4\%}; cost of risk $\sim$\textbf{114bps}. Capital return is the standout: on track for \textbf{$\geq$\texteuro10bn of buybacks across 2025--26} on top of a rising dividend (FY25 \texteuro0.24/share, +14\%). Market cap $\sim$\texteuro150bn, $\sim$10x P/E. \textbf{Thesis driver:} the $\sim$\$12.3bn \textbf{Webster Financial} acquisition (close expected H2 2026) builds a top-10 US bank, while the Poland exit ($\sim$\texteuro1.9bn gain) funds simplification and the buyback (\href{https://investors.websterbank.com/News--Events/news-releases/news-details/2026/Webster-Financial-Corporation-Enters-Into-Merger-Agreement-With-Banco-Santander-S-A--for-12-3-Billion/default.aspx}{Webster IR}).

\subsection*{ING Group (INGA)}
Benelux-anchored retail/wholesale bank built on a large, low-cost deposit franchise ($\sim$45\% Dutch retail share; ING-DiBa the leading German direct bank). Q1 net result \textbf{\texteuro1,556m (+6.9\%)}; commercial NII \textbf{\texteuro4,060m (+12\%)}; fee income \texteuro1,236m (+13\%); RoTE \textbf{13.6\%} (4-quarter rolling 13.9\%) (\href{https://www.globenewswire.com/news-release/2026/04/30/3284411/0/en/}{ING 1Q26}). CET1 \textbf{13.0\%} vs $\sim$13\% target; cost of risk \textbf{19bps}. Capital return: new \textbf{\texteuro1.0bn buyback} announced in Q1, continuing the policy of distributing capital above the $\sim$13\% CET1 floor. Most expensive of the eight at $\sim$1.48x P/TBV; market cap $\sim$\$88--89bn. \textbf{Thesis driver:} premium, sticky deposit franchise delivering sector-leading RoTE with aggressive, repeatable buybacks; metric shift from RoE to RoTE signals confidence in capital efficiency.

\subsection*{UniCredit (UCG)}
Pan-European commercial bank (Italy, Germany's HVB, Austria, CEE) with a strong fee/wealth/trading mix, run for best-in-class profitability under Andrea Orcel. Q1 net profit a record \textbf{\texteuro3.2bn (+16\%)}; NII \textbf{\texteuro3.59bn ($-$2\% YoY)}, the decline offset by fees and trading; RoTE \textbf{25.8\%}, the highest of the eight (\href{https://www.unicreditgroup.eu/en/press-media/press-releases-price-sensitive/2026/may/unicredit--1q26-group-results---unlimited-off-to-a-flying-start.html}{UniCredit 1Q26}). CET1 \textbf{14.2\%} (14.8\% pro-forma); cost of risk a strikingly low $\sim$\textbf{8bps}. FY26 net-profit target lifted to \textbf{$\geq$\texteuro11bn}. Trades $\sim$1.88x tangible book, $\sim$8.5x forward P/E. \textbf{Thesis driver:} M\&A optionality dominates --- the \textbf{Commerzbank} stake crossed $\sim$34\% as of 2 June 2026 (\href{https://www.bloomberg.com/news/articles/2026-06-02/unicredit-says-commerzbank-offer-set-to-push-direct-stake-to-34}{Bloomberg}), the \textbf{Generali} stake rose to $\sim$8.7\%, and the \textbf{Banco BPM} bid was withdrawn after Rome's ``golden power'' conditions.

\subsection*{Intesa Sanpaolo (ISP)}
Italy's largest bank; an integrated domestic retail/commercial model with a fast-growing, fee-rich Wealth Management, Protection \& Advisory division. Q1 net profit a record \textbf{\texteuro2.8bn (+6\%)}, annualised ROE $\sim$\textbf{21\%}; NII \textbf{\texteuro3.6bn} (roughly flat despite lower rates, protected by core-deposit hedging) (\href{https://group.intesasanpaolo.com/en/newsroom/all-news/news/2026/intesa-sanpaolo-1q26-results}{Intesa 1Q26}). CET1 \textbf{>13.0\%}; cost of risk a very low $\sim$\textbf{16bps}. FY26 net income confirmed $\sim$\textbf{\texteuro10bn}. Capital return is the headline: $\sim$\textbf{\texteuro9.4bn planned for 2026} (final + interim dividends plus a \textbf{\texteuro2.3bn buyback launching July 2026}), a $\sim$95\% payout and $\sim$7.5\% yield. Now among Europe's largest banks by market cap ($\sim$\texteuro100bn+). \textbf{Thesis driver:} best-in-class shareholder yield funded by a capital-light fee engine and ultra-low cost of risk; chief risk is ECB-cut pressure on NII.

\subsection*{Deutsche Bank (DBK)}
Germany-anchored ``Global Hausbank'' with four divisions (Corporate Bank, Investment Bank, Private Bank, DWS), all delivering $\geq$13\% RoTE in Q1. Record post-tax profit \textbf{\texteuro2.2bn (+8\%)} on revenue \textbf{\texteuro8.7bn (+2\%)}; RoTE \textbf{12.7\%}, up from 11.9\% (\href{https://www.db.com/news/detail/20260429-deutsche-bank-reports-first-quarter-2026-results}{Deutsche Bank 1Q26}). CET1 \textbf{13.8\%} (13.5--14.0\% range); cost of risk guided toward $\sim$30bps; banking-book \textbf{NII guided to $\sim$\texteuro14bn for 2026} on structural-hedge roll-over. Capital return: $\sim$60\% payout target and a \textbf{\texteuro1bn buyback} launched Feb 2026. Re-rated to $\sim$0.77x P/TBV, market cap $\sim$\texteuro60bn. \textbf{Thesis driver:} delivered its 2025 targets (record \texteuro9.7bn PBT, 10.3\% RoTE) and launched the next ``Scaling the Global Hausbank'' phase targeting \textbf{>13\% RoTE by 2028} --- Q1 is the proof-point quarter.

\subsection*{BBVA (BBVA)}
Retail/commercial bank weighted to high-growth EM, led by Mexico, plus Spain, Turkey (Garanti) and South America. Q1 net profit \textbf{\texteuro2.99bn (+10.8\%)}, beating consensus; NII \textbf{\texteuro7.54bn (+20\% at constant FX)} on $\sim$17\% lending growth; RoTE \textbf{21.7\%} (\href{https://www.bbva.com/en/economy-and-finance/earnings-1q2026/}{BBVA 1Q26}). CET1 \textbf{12.83\%} (above the 11.5--12.0\% target); cost of risk $\sim$\textbf{154bps}, structurally higher given the EM/consumer mix. Capital return: launching the \textbf{final \texteuro1.46bn tranche} of a $\sim$\texteuro4.6bn post-Sabadell buyback. \textbf{Thesis driver:} the \textbf{\texteuro17bn hostile bid for Banco Sabadell failed} in October 2025 (only $\sim$25.5\% tendered); the stock then rallied $\sim$34\% as BBVA pivoted to a standalone capital-return story (\href{https://www.euronews.com/business/2025/10/17/bbva-fails-in-17bn-takeover-battle-for-smaller-spanish-rival-sabadell}{euronews}).

\subsection*{Soci\'et\'e G\'en\'erale (GLE)}
French universal bank --- domestic retail (incl.\ BoursoBank), Global Banking \& Investor Solutions (a leader in equity derivatives), and the Ayvens mobility/leasing arm. Q1 net income \textbf{\texteuro1,696m (+5.5\%)} on NBI \textbf{\texteuro7,106m (+0.3\%)}; French-retail \textbf{NII +12\% YoY} as legacy hedge drag rolls off; reported RoTE \textbf{11.7\%} (12.7\% adjusted) (\href{https://www.globenewswire.com/news-release/2026/04/30/3284407/0/en/societe-generale-first-quarter-2026-earnings.html}{SocGen 1Q26}). CET1 a robust \textbf{13.5\%}; cost of risk \textbf{25bps}. Capital return: 2025 cash dividend \texteuro1.61/share plus a \texteuro1,462m buyback completed March 2026. Cheapest of the eight at $\sim$0.6--0.8x P/TBV, market cap $\sim$\texteuro52bn. \textbf{Thesis driver:} self-help --- opex $-$6\% YoY, disposals, and the French-retail NII recovery are lifting RoTE off a deep-discount base; SocGen raised its 2026 profitability target.

\subsection*{Cross-issuer comparison (Q1 2026)}

\noindent\footnotesize
\begin{tabularx}{\textwidth}{@{}l p{1.7cm} p{2cm} p{0.9cm} p{1cm} p{1.4cm} p{1.7cm} p{1.6cm} p{1.1cm}@{}}
\toprule
\textbf{Issuer} & \textbf{Net profit (YoY)} & \textbf{NII trend} & \textbf{CET1} & \textbf{CoR} & \textbf{RoTE} & \textbf{'25--'26 buyback} & \textbf{Mkt cap} & \textbf{P/TBV} \\
\midrule
BNP Paribas       & \texteuro3,217m (+9\%)              & comm. bank $\sim$+7\%   & 12.8\%   & 39bps    & 12\% FY26 tgt      & \texteuro1.15bn done       & $\sim$\texteuro102bn        & 0.9--1.0x \\
Santander         & \texteuro3.56bn (+12\%)             & +4\%                    & 14.4\%   & $\sim$114bps & 15.2\%         & $\geq$\texteuro10bn '25--'26 & $\sim$\texteuro150bn*       & n/a \\
ING               & \texteuro1,556m (+6.9\%)            & comm. NII +12\%         & 13.0\%   & 19bps    & 13.6\%             & \texteuro1.0bn new         & $\sim$\$88--89bn            & $\sim$1.48x \\
UniCredit         & \texteuro3.2bn (+16\%)              & $-$2\%                  & 14.2\%   & $\sim$8bps & 25.8\%           & new tranche                & $\sim$\texteuro90--100bn*   & $\sim$1.88x \\
Intesa Sanpaolo   & \texteuro2.8bn (+6\%)               & $\sim$flat              & >13.0\%  & $\sim$16bps & $\sim$21\% ROE  & \texteuro2.3bn Jul '26     & $\sim$\texteuro100bn+*      & n/a \\
Deutsche Bank     & \texteuro2.2bn (+8\%)               & $\sim$\texteuro14bn FY26 & 13.8\%   & $\sim$30bps & 12.7\%         & \texteuro1.0bn             & $\sim$\texteuro60bn         & $\sim$0.77x \\
BBVA              & \texteuro2.99bn (+10.8\%)           & +20\% (const. FX)       & 12.83\%  & $\sim$154bps & 21.7\%        & \texteuro4.6bn prog.       & $\sim$\texteuro115--125bn*  & n/a \\
Soci\'et\'e G\'en\'erale & \texteuro1,696m (+5.5\%)     & Fr. retail +12\%        & 13.5\%   & 25bps    & 11.7\%             & \texteuro1,462m done       & $\sim$\texteuro52bn         & 0.6--0.8x \\
\bottomrule
\end{tabularx}
\normalsize

\smallskip
\noindent\textit{\footnotesize * Market caps are source-conflicted and approximate (see Section 6). RoTE figures are not strictly comparable --- Intesa reports ROE, and BNP/SocGen FY-target vs reported bases differ.}

% =============================================================================
\section{Q1 2026 cross-issuer themes}

\subsection*{NII sensitivity to the ECB}
With the deposit rate parked at 2.00\% and the easing cycle apparently over, the NII story has bifurcated. The names that suffered most from cuts are now seeing the drag roll off: \textbf{SocGen's French-retail NII rose +12\%} and \textbf{BNP's commercial bank accelerated to $\sim$+7\%} as legacy hedges reset and the curve steepened. The EM-geared books were insulated all along --- \textbf{BBVA's NII grew +20\% at constant FX} on Mexican volume, and \textbf{Santander's rose +4\%}. The most rate-sensitive domestically are the Italians, but both have managed it: \textbf{Intesa held NII roughly flat} via core-deposit hedging, and \textbf{UniCredit absorbed a $-$2\% NII decline} with fees and trading. \textbf{Deutsche guided banking-book NII up to $\sim$\texteuro14bn} on structural-hedge roll-over. Net read: the sector's NII headwind has largely passed, and a steeper curve is now a tailwind into H2 2026. No bank disclosed a clean euros-per-25bp sensitivity, so this is directional rather than an elasticity.

\subsection*{Capital return}
This is the sector's defining feature in 2026 and the clearest differentiator. \textbf{Santander} leads on absolute scale (\textbf{$\geq$\texteuro10bn of buybacks across 2025--26}), \textbf{Intesa} on yield (\textbf{$\sim$\texteuro9.4bn / $\sim$95\% payout / $\sim$7.5\%} for 2026, with a \texteuro2.3bn buyback from July), and \textbf{BBVA} on the post-deal pivot (the final \texteuro1.46bn of a $\sim$\texteuro4.6bn programme after the Sabadell bid collapsed). The universal banks run smaller but steady programmes --- BNP (\texteuro1.15bn, done), SocGen (\texteuro1,462m, done), Deutsche (\texteuro1bn, $\sim$60\% payout), ING (a fresh \texteuro1bn, distributing everything above $\sim$13\% CET1). Sector dividend yields sit around \textbf{4.8\%}, and index-wide STOXX 600 dividends are forecast +4\% in 2026 (\href{https://www.euronews.com/business/2025/12/13/european-banks-best-record-whats-the-outlook-for-2026}{euronews}). Strong CET1 ratios (12.8\%--14.4\% across the group) underwrite it.

\subsection*{Cost-of-risk normalisation}
Cost of risk remains \textbf{below cycle average} across the group, with a moderate pick-up expected from H2 2026 (\href{https://www.scoperatings.com/ratings-and-research/research/EN/179762}{Scope}). The spread is wide and mix-driven: UniCredit ($\sim$8bps), Intesa ($\sim$16bps) and ING (19bps) sit at the benign end; the universal banks cluster at 25--39bps (SocGen 25, Deutsche $\sim$30, BNP 39); and the EM names run structurally higher (Santander $\sim$114bps, BBVA $\sim$154bps) by design, compensated by far higher margins. Nothing in Q1 points to credit deterioration --- provisions are still being released or run below guidance at several names.

\subsection*{CRE exposure}
Commercial real estate remains the watch-item but not the scare-item. The ECB's May 2026 Financial Stability Review judges \textbf{systemic CRE exposure contained and unlikely to threaten euro-area bank solvency}, while flagging residual pockets of stress (CRE and post-tariff auto lending) (\href{https://www.ecb.europa.eu/press/financial-stability-publications/fsr/html/ecb.fsr202605~50566915a7.en.html}{ECB FSR}). The ECB's SSM has nonetheless put CRE explicitly under its 2026--2028 supervisory microscope. None of the eight flagged a material CRE provisioning event in Q1, but disclosure granularity varies --- a follow-up item below.

% =============================================================================
\section{Regulatory backdrop}

\textbf{Basel III endgame (CRR3/CRD6).} CRR3 has applied since \textbf{1 January 2025}; CRD6 transposition into national law is due in early 2026 (\href{https://www.pwc.com/gx/en/services/audit-assurance/risk-assurance/ssm-eba-office/eu-publishes-crr3-crd6-official-journal.html}{PwC}). The market-risk piece (FRTB) keeps slipping: after a first delay, the Commission in June 2025 pushed application from Jan 2026 to \textbf{Jan 2027} to protect EU competitiveness, and on \textbf{4 June 2026} adopted a further delegated act adding a targeted multiplier and operational relief, phased Jan 2027--Jan 2030 (\href{https://finance.ec.europa.eu/news/eu-temporarily-amends-prudential-rules-banks-market-risk-2026-06-08_en}{EC, Jun 2026}). Net effect: the capital headwind from Basel finalisation is being smoothed and deferred --- supportive of the capital-return theme above.

\textbf{EBA EU-wide stress test 2025 (results Aug 2025).} Across 64 banks ($\sim$75\% of EU assets), the adverse scenario depleted CET1 by \textbf{\texteuro229bn, taking the aggregate ratio from 15.76\% to 12.06\% ($-$370bps)} --- milder than the $-$479bps of the 2023 test --- with \textbf{no bank breaching its TSCR} and \texteuro554bn of CET1 held above requirements (\href{https://www.eba.europa.eu/publications-and-media/publications/2025-eu-wide-stress-test-results}{EBA}). This was the EBA EU-wide exercise (the ECB ran its SSM component within it); there was no separate standalone SSM 2025 test result.

\textbf{SSM supervisory priorities 2026--2028 (Nov 2025).} Two themes: financial/geopolitical resilience (tighter scrutiny of new lending, SMEs and \textbf{CRE} in focus) and operational/ICT resilience (cyber, third-party and AI-reliance risk) (\href{https://www.bankingsupervision.europa.eu/framework/priorities/html/ssm.supervisory_priorities202511.en.html}{ECB Banking Supervision}).

\textbf{MREL.} Effectively a solved problem for the named issuers: the EBA's Q2 2025 dashboard shows an aggregate shortfall of just \textbf{\texteuro0.3bn (<0.01\% of TREA)}, with all banks past their final-target date compliant; the residual sits with names still in transition. $\sim$\texteuro221bn of instruments become ineligible by June 2026, a rollover (not a capital) need (\href{https://www.eba.europa.eu/publications-and-media/press-releases/eba-publishes-its-q2-2025-dashboard-minimum-requirement-own-funds-and-eligible-liabilities}{EBA Q2 2025}).

\textbf{Banking union / CMDI / M\&A friction.} The CMDI crisis-management reform was finalised --- Council/Parliament deal June 2025, Parliament approval March 2026 (\href{https://www.consilium.europa.eu/en/press/press-releases/2025/06/25/bank-resolution-council-and-parliament-strike-deal-to-strengthen-the-eu-crisis-management-framework/}{Council}). But national-level friction is now the dominant M\&A story: Spain's conditions helped sink \textbf{BBVA--Sabadell} (failed Oct 2025; the Commission opened an infringement case against Spain), Rome's ``golden power'' killed \textbf{UniCredit--Banco BPM}, and Berlin is resisting \textbf{UniCredit--Commerzbank} even as UniCredit's stake crosses $\sim$34\% (closing not expected before 2027). The lesson for the sector: cross-border consolidation logic is strong, but the binding constraint is political, not prudential.

% =============================================================================
\section{Items flagged for analyst follow-up}

\begin{enumerate}
\item \textbf{UniCredit--Commerzbank --- track the path past 34\%.} The acceptance period runs to $\sim$3 July 2026 and closing is a 2027 event pending approvals against Berlin opposition. Watch for the EU's stance on Germany's resistance (parallel to the Spain infringement case) and any move by UniCredit to consolidate vs.\ hold as a financial stake --- it changes the pro-forma capital and Generali optionality materially.

\item \textbf{Intesa's \texteuro2.3bn buyback execution (July 2026 start).} Confirm the programme launches on schedule and whether the $\sim$95\% payout is reaffirmed at H1 results --- it is the single biggest support for the ISP thesis and the most exposed to any NII disappointment if the ECB surprises hawkish.

\item \textbf{ECB 11 June 2026 decision and guidance.} The market has flipped from pricing cuts to debating a hike. A hawkish surprise re-rates the NII-geared names (SocGen, BNP, the Italians) positively and pressures the EM-funding cost of Santander/BBVA. Read the staff projections for the 2026--27 rate path the banks are underwriting.

\item \textbf{Santander--Webster close (H2 2026).} Watch for the regulatory approval timeline and any revision to accretion/integration guidance; a delayed or repriced close would dent the US-scale thesis and the buyback cadence funded partly by the Poland gain.

\item \textbf{CRE disclosure granularity at H1 2026.} The ECB calls systemic CRE risk contained but has it under supervisory focus. Push each issuer for office-vs-other CRE splits, LTVs and watch-list migration at the half-year --- the gap between ``contained at system level'' and name-level concentration is where surprises live.
\end{enumerate}

% =============================================================================
\section{What I assumed / what I could not verify}

\begin{itemize}
\item \textbf{Market caps and P/TBV are approximate and, in places, source-conflicted.} Values come from third-party aggregators (companiesmarketcap, gurufocus, Yahoo, CNBC) dated late-May to 7-June 2026, not company filings. In particular: Santander shows $\sim$\texteuro150bn on one aggregator while a separate source calls \textbf{Intesa} ``Europe's largest bank by market cap'' ($\sim$\texteuro100bn+) --- these cannot both be precisely right, so treat the league-table ordering as indicative and re-pull on a live terminal. BBVA's market cap ranged from an implausible $\sim$\texteuro53bn print to a $\sim$\texteuro114bn figure; $\sim$\texteuro115--125bn is the likelier number post-rally. P/TBV was only cleanly sourced for ING ($\sim$1.48x), UniCredit ($\sim$1.88x), Deutsche ($\sim$0.77x) and SocGen ($\sim$0.6--0.8x); BNP is ``near book,'' and Santander/BBVA P/TBV were not isolated (P/E used instead).

\item \textbf{No bank disclosed a clean NII sensitivity to a 25bp ECB move.} All rate-sensitivity commentary in Section 3 is directional, inferred from divisional NII trends and management hedging commentary --- not a disclosed euros-per-basis-point elasticity. The closest hard assumption is BNP's Belgian-retail plan (ECB $\sim$2.5\%$\rightarrow$2.0\% by end-2027; long-end $\sim$2.8\%).

\item \textbf{FY2026 group NII guidance is uneven across the group.} Deutsche ($\sim$\texteuro14bn) and ING (total income $\sim$\texteuro24bn) guide clearly; BNP, SocGen, Santander and BBVA guide on revenue/RoTE rather than a single consolidated NII line, so their NII ``guidance'' here is a proxy from divisional trends.

\item \textbf{The ECB 11 June 2026 outcome is unknown at time of writing.} The ``market leaning toward a hike'' read is well-corroborated directionally (CNBC, JPM commentary) but the precise market-implied probability seen ($\sim$99\%) came from a single commercial source and should not be quoted as fact.

\item \textbf{A precise YTD 2026 return and a current P/TBV for the STOXX 600 Banks index were not isolated} --- the index is qualitatively down from its 3 Feb 2026 high of $\sim$385 toward $\sim$370, but no exact YTD figure or dated June P/TBV print was obtained.

\item \textbf{Deutsche Bank's Q1 profit is cited two ways} --- \texteuro2.2bn post-tax (the headline used here) vs \texteuro1.9bn attributable to shareholders. Both are correct; they measure different lines.

\item \textbf{Intesa reports ROE ($\sim$21\%), not a separate RoTE}, so its return figure is not strictly comparable to the RoTE figures shown for the other seven.
\end{itemize}

% =============================================================================
\section*{Sources}

\footnotesize
\textbf{Company results:}
\href{https://invest.bnpparibas/en/document/1q26-pr}{BNP 1Q26} $\cdot$
\href{https://www.santander.com/en/press-room/press-releases/2026/04/q1-2026-santander-bank-results}{Santander Q1 2026} $\cdot$
\href{https://www.globenewswire.com/news-release/2026/04/30/3284411/0/en/}{ING 1Q26} $\cdot$
\href{https://www.unicreditgroup.eu/en/press-media/press-releases-price-sensitive/2026/may/unicredit--1q26-group-results---unlimited-off-to-a-flying-start.html}{UniCredit 1Q26} $\cdot$
\href{https://group.intesasanpaolo.com/en/newsroom/all-news/news/2026/intesa-sanpaolo-1q26-results}{Intesa 1Q26} $\cdot$
\href{https://www.db.com/news/detail/20260429-deutsche-bank-reports-first-quarter-2026-results}{Deutsche Bank 1Q26} $\cdot$
\href{https://www.bbva.com/en/economy-and-finance/earnings-1q2026/}{BBVA 1Q26} $\cdot$
\href{https://www.globenewswire.com/news-release/2026/04/30/3284407/0/en/societe-generale-first-quarter-2026-earnings.html}{SocGen 1Q26}

\textbf{Macro \& regulatory:}
\href{https://www.ecb.europa.eu/stats/policy_and_exchange_rates/key_ecb_interest_rates/html/index.en.html}{ECB key rates} $\cdot$
\href{https://www.ecb.europa.eu/press/financial-stability-publications/fsr/html/ecb.fsr202605~50566915a7.en.html}{ECB FSR May 2026} $\cdot$
\href{https://www.eba.europa.eu/publications-and-media/publications/2025-eu-wide-stress-test-results}{EBA 2025 stress test} $\cdot$
\href{https://www.eba.europa.eu/publications-and-media/press-releases/eba-publishes-its-q2-2025-dashboard-minimum-requirement-own-funds-and-eligible-liabilities}{EBA Q2 2025 MREL} $\cdot$
\href{https://finance.ec.europa.eu/news/eu-temporarily-amends-prudential-rules-banks-market-risk-2026-06-08_en}{EC FRTB Jun 2026} $\cdot$
\href{https://www.bankingsupervision.europa.eu/framework/priorities/html/ssm.supervisory_priorities202511.en.html}{SSM priorities 2026--28} $\cdot$
\href{https://stoxx.com/european-bank-stocks-have-record-year-lifting-sector-stoxx-etf-assets-above-eur-13bn/}{STOXX banks 2025}
\normalsize

\vfill
\noindent{\footnotesize\color{darkgray}\textit{Generated via personal Claude Cowork instance on public data, for office demonstration only. Not produced using office tooling or office data.}}

\end{document}
